\subsection{带标架的半单李代数}

命题 1. 设 \((\mathfrak{g},\mathfrak{h})\) 是分裂半单李代数，R 是它的根系，B 是 R 的基，\((n(\alpha ,\beta))_{\alpha ,\beta \in \mathrm{B}}\) 是对应的嘉当矩阵。对一切 \(\alpha \in \mathrm{B}\) ，设 \(X_{\alpha}\in \mathfrak{g}^{\alpha}\) ，\(X_{-\alpha}\in \mathfrak{g}^{-\alpha}\) 。则对 \(\alpha ,\beta \in \mathrm{B}\) ，

\[
[ H _ {\alpha}, H _ {\beta} ] = 0\tag{1}
\]

\[
[ H _ {\alpha}, X _ {\beta} ] = n (\beta , \alpha) X _ {\beta}\tag{2}
\]

\[
[ H _ {\alpha}, X _ {- \beta} ] = - n (\beta , \alpha) X _ {- \beta}\tag{3}
\]

\[
[ X _ {- \alpha}, X _ {\beta} ] = 0 \quad \text { if } \alpha \neq \beta\tag{4}
\]

\[
(\operatorname{ad} X _ {\alpha}) ^ {1 - n (\beta , \alpha)} X _ {\beta} = 0 \quad \text { if } \alpha \neq \beta\tag{5}
\]

\[
(\operatorname{ad} X _ {- \alpha}) ^ {1 - n (\beta , \alpha)} X _ {- \beta} = 0 \quad \text { if } \alpha \neq \beta .\tag{6}
\]

族 \((H_{\alpha})_{\alpha \in \mathrm{B}}\) 是 \(\mathfrak{h}\) 的一个基。若 \(X_{\alpha} \neq 0\) 且 \(X_{-\alpha} \neq 0\) 对一切 \(\alpha \in \mathrm{B}\) 成立，则李代数 \(\mathfrak{g}\) 由诸 \(X_{\alpha}\) 与诸 \(X_{-\alpha}\) （ \(\alpha \in \mathrm{B}\) ）生成。

（回忆，若 \(\alpha, \beta \in \mathrm{B}\) 且 \(\alpha \neq \beta\) ，则 \(n(\beta, \alpha)\) 是一个 \(\leq 0\) 的整数，故公式 (5) 与 (6) 有意义。）

公式 (1)、(2) 与 (3) 是明显的。若 \(\alpha \neq \beta\) ，则 \(\beta - \alpha\) 不是根，因为 \(R\) 的每个元素都是 \(B\) 的元素的具整数系数的线性组合，且这些系数全同号（第六章 §1 第 6 节 定理 3）。这就证明了 (4)。鉴于第六章 §1 第 3 节 命题 9，这还证明了 \(\alpha\) -链（由 \(\beta\) 定义的）是

\[
\{\beta , \beta + \alpha , \dots , \beta - n (\beta , \alpha) \alpha \};
\]

从而 \(\beta + (1 - n(\beta, \alpha))\alpha \notin \mathrm{R}\) ，这就证明了 (5)。等式 (6) 用类似的方式建立。族 \((H_{\alpha})_{\alpha \in \mathrm{B}}\) 是 \(\mathrm{R}^{\vee}\) 的、从而是 \(\mathfrak{h}\) 的基。若 \(X_{\alpha} \neq 0\) 且 \(X_{-\alpha} \neq 0\) 对一切 \(\alpha \in \mathrm{B}\) 成立，则 \([X_{\alpha}, X_{-\alpha}] = \lambda_{\alpha}H_{\alpha}\) ，其中 \(\lambda_{\alpha} \neq 0\) ，故最后的断言由 §3 第 3 节 命题 9 (iii) 得出。

定义 1. 设 \((\mathfrak{g},\mathfrak{h})\) 是分裂半单李代数，R 是它的根系。\((\mathfrak{g},\mathfrak{h})\) 的一个标架是指一个偶 \((\mathrm{B},(X_{\alpha})_{\alpha \in \mathrm{B}})\) ，其中 B 是 R 的一个基，并且对一切 \(\alpha \in \mathrm{B}\) ，\(X_{\alpha}\) 是 \(\mathfrak{g}^{\alpha}\) 的一个非零元素。带标架的半单李代数是指一个序列 \((\mathfrak{g},\mathfrak{h},\mathrm{B},(X_{\alpha})_{\alpha \in \mathrm{B}})\) ，其中 \((\mathfrak{g},\mathfrak{h})\) 是分裂半单李代数，且 \((\mathrm{B},(X_{\alpha})_{\alpha \in \mathrm{B}})\) 是 \((\mathfrak{g},\mathfrak{h})\) 的一个标架。

\(\mathfrak{g}\) 的一个标架是指 \((\mathfrak{g},\mathfrak{h})\) 的一个标架，其中 \(\mathfrak{h}\) 是 \(\mathfrak{g}\) 的一个分裂嘉当子代数。

设 \(a_1 = (\mathfrak{g}_1, \mathfrak{h}_1, \mathrm{B}_1, (X_\alpha^1)_{\alpha \in \mathrm{B}_1})\) 与 \(a_2 = (\mathfrak{g}_2, \mathfrak{h}_2, \mathrm{B}_2, (X_\alpha^2)_{\alpha \in \mathrm{B}_2})\) 是带标架的半单李代数。从 \(a_1\) 到 \(a_2\) 的一个同构是指一个同构 \(\varphi\) （从 \(\mathfrak{g}_1\) 到 \(\mathfrak{g}_2\) ），它把 \(\mathfrak{h}_1\) 映到 \(\mathfrak{h}_2\) ，把 \(\mathrm{B}_1\) 映到 \(\mathrm{B}_2\) ，并且把 \(X_\alpha^1\) 映到 \(X_{\psi \alpha}^2\) 对一切 \(\alpha \in \mathrm{B}_1\) 成立（其中 \(\psi\) 是 \(\varphi | \mathfrak{h}_1\) 的反步映射）。这时，称 \(\varphi\) 把标架 \((\mathrm{B}_1, (X_\alpha^1)_{\alpha \in \mathrm{B}_1})\) 变换为标架 \((\mathrm{B}_2, (X_\alpha^2)_{\alpha \in \mathrm{B}_2})\) 。

若 \((\mathrm{B},(X_{\alpha})_{\alpha\in\mathrm{B}})\) 是 \((\mathfrak{g},\mathfrak{h})\) 的一个标架，则对一切 \(\alpha\in B\) ，存在唯一的元素 \(X_{-\alpha}\) （\(g^{-\alpha}\) 中的）使 \([X_{\alpha},X_{-\alpha}]=-H_{\alpha}\) （§2 第 2 节 定理 1 (iv)）。族 \((X_{\alpha})_{\alpha\in\mathrm{B}\cup(-\mathrm{B})}\) 称为由该标架定义的生成族（参见命题 1）。它也是由标架 \((-\mathrm{B},(X_{\alpha})_{\alpha\in-\mathrm{B}})\) 定义的生成族。对一切 \(\alpha\in\mathrm{B}\cup(-\mathrm{B})\) ，设 \(t_{\alpha}\in k^{*}\) ，并设 \(t_{\alpha}t_{-\alpha}=1\) 对一切 \(\alpha\in B\) 成立。则 \((t_{\alpha}X_{\alpha})_{\alpha\in\mathrm{B}\cup(-\mathrm{B})}\) 是由标架 \((\mathrm{B},(t_{\alpha}X_{\alpha})_{\alpha\in\mathrm{B}})\) 定义的生成族。

